Consider an equilibrium position of an autonomous Hamiltonian system, assumed stable in linear approximation. Let $\omega_1, \ldots, \omega_n$ be the $n$ characteristic frequencies. Assume no low-order resonance relations:
$$k_1 \omega_1 + \cdots + k_n \omega_n = 0 \quad \text{with integers } k_i \text{ such that } 0 < \sum |k_i| \leq 4.$$

Then the Hamiltonian can be reduced to Birkhoff normal form:
$$H = H_0(\tau) + \cdots, \qquad H_0(\tau) = \sum \omega_k \tau_k + \frac{1}{2} \sum \omega_{kl} \tau_k \tau_l,$$
where $\tau_k = (p_k^2 + q_k^2)/2$ and the dots denote terms of degree higher than four in the distance from equilibrium.

The nondegeneracy condition
$$\det |\omega_{kl}| \neq 0$$
guarantees the existence of a set of invariant tori of almost full measure in a sufficiently small neighborhood of the equilibrium position.

The isoenergetic nondegeneracy condition
$$\det \begin{vmatrix} \omega_{kl} & \omega_k \\ \omega_l & 0 \end{vmatrix} \neq 0$$
guarantees such a set on every energy level set sufficiently close to the critical point.

In the case $n = 2$, isoenergetic nondegeneracy is satisfied if the quadratic part of $H_0$ is not divisible by the linear part, and it guarantees Liapunov stability of the equilibrium position.
